\documentclass[main.tex]{subfiles}
\begin{document}

\phantomsection
\addcontentsline{toc}{chapter}{Lời mở đầu}
\chapter*{Lời mở đầu}
    
Trong bối cảnh nền kinh tế thị trường và sự phát triển mạnh mẽ của chuỗi cung ứng toàn cầu, việc tối ưu hóa chi phí logistics đóng vai trò sống còn đối với mọi doanh nghiệp sản xuất và phân phối. Một trong những mô hình toán học kinh điển và thành công nhất được ứng dụng để giải quyết vấn đề này là bài toán vận tải. 
    
Về mặt lịch sử, bài toán vận tải lần đầu tiên được định hình và phát biểu rõ ràng dưới dạng toán học bởi nhà toán học người Mỹ Frank Lauren Hitchcock vào năm 1941 \cite{hitchcock1941}. Hitchcock đã đặt nền móng cho việc mô hình hóa bài toán phân phối một sản phẩm từ nhiều nguồn cung đến nhiều điểm cầu sao cho tổng chi phí là nhỏ nhất. Tiếp đó vào năm 1947, nhà kinh tế học Tjalling C. Koopmans công bố công trình nghiên cứu độc lập của mình \cite{koopmans1949}, trong đó ông nghiên cứu sâu sắc bài toán này dưới góc độ kinh tế và logistics hàng hải, đồng thời đưa ra những khái niệm đầu tiên về hệ thống thế vị tương đương với chi phí biên. Sự đóng góp của Koopmans lớn đến mức bài toán này thường được gọi là bài toán Hitchcock-Koopmans.
    
Sự đột phá về phương pháp giải chỉ thực sự đến khi George B. Dantzig (1951) \cite{dantzig1951} - cha đẻ của Quy hoạch tuyến tính - áp dụng thuật toán đơn hình vào bài toán vận tải. Trong chuyên khảo số 13 của Cowles Commission, Dantzig đã chỉ ra rằng cấu trúc ma trận thưa và đặc biệt của bài toán vận tải cho phép đơn giản hóa thuật toán đơn hình thành một dạng cực kỳ hiệu quả, tính toán trực tiếp trên bảng mà không cần nghịch đảo ma trận. Phương pháp này chính là thuật toán thế vị hay thuật toán đơn hình trên mạng.

Nội dung lý thuyết trọng tâm của đồ án được tham khảo và chuẩn hóa từ giáo trình Các phương pháp Tối ưu của tác giả Nguyễn Thị Bạch Kim \cite{bachkim}, kết hợp với các tài liệu nghiên cứu gốc để đảm bảo tính hàn lâm và chặt chẽ của toán học. Mặc dù đã nỗ lực hoàn thiện, đồ án khó tránh khỏi những thiếu sót, rất mong nhận được sự góp ý của thầy giáo hướng dẫn để đồ án được hoàn thiện hơn.

\end{document}